% La Rutina Diaria -- Worksheet 4 (reflexive verbs)
% ---------------------------------------------------------------------------
% Cleaned from inputs/spanish/spanish-rutina-diaria-DRAFT.docx by the handout-formatting
% skill. THIS IS THE SINGLE SOURCE. The body is written once; the audience is a
% compile option:
%   tectonic rutina-diaria.tex      -> student handout (blanks, no answers)
%   tectonic rutina-diaria-key.tex  -> answer key      (answers shown)
% The -key file is a two-line wrapper that sets \audience and \input this file
% -- never edit the body to make a key. (\usepackage[teacher]{handout} also
% works, adding the teacher notes, if you ever want that build.)
% ---------------------------------------------------------------------------
\documentclass[11pt]{article}
\usepackage{housestyle}
\providecommand{\audience}{}      % empty = student; wrappers override with key/teacher
\usepackage[\audience]{handout}

\hypersetup{pdftitle={La Rutina Diaria -- Worksheet 4}, pdflang={es-ES}}

\begin{document}

\handouttitle{La Rutina Diaria}{Worksheet 4}{}

\calloutbox{Objetivos de Hoy}{%
    \begin{itemize}[leftmargin=*]
        \item Reconocer y usar los verbos reflexivos para hablar de la rutina diaria.
        \item Colocar correctamente el pronombre reflexivo (\textit{me, te, se, nos, os, se}).
        \item Conjugar verbos reflexivos en el presente.
        \item Leer y comprender un texto corto sobre un día normal.
    \end{itemize}%
}

\bigskip

\section{Lectura}

\begin{passage}
Me llamo Lucía y todos los días me levanto a las seis y media de la mañana.
Primero me ducho y me lavo el pelo; después me visto y desayuno café con leche
y pan tostado. Mi hermano se despierta más tarde porque se acuesta muy tarde, y
siempre se queja de que tiene sueño. Nosotros nos cepillamos los dientes antes
de salir. Por la noche, mi familia y yo nos sentamos juntos para cenar y nos
relajamos un poco antes de acostarnos.
\end{passage}

\section{Vocabulario}

\begin{vocabtable}
    \vocabentry{despertarse}{to wake up}
    \vocabentry{levantarse}{to get up}
    \vocabentry{ducharse}{to shower}
    \vocabentry{lavarse}{to wash (oneself)}
    \vocabentry{vestirse}{to get dressed}
    \vocabentry{acostarse}{to go to bed}
    \vocabentry{cepillarse los dientes}{to brush one's teeth}
    \vocabentry{sentarse}{to sit down}
    \vocabentry{quejarse (de)}{to complain (about)}
    \vocabentry{relajarse}{to relax}
\end{vocabtable}

\section{Conjugación: \textit{levantarse} (presente)}

\begin{center}\renewcommand{\arraystretch}{1.3}
\begin{tabular}{@{} l l l @{}}
    \toprule
    \textbf{Pronombre} & \textbf{Forma} & \textbf{English} \\
    \midrule
    yo                & me levanto      & I get up \\
    tú                & te levantas     & you get up \\
    él / ella / Ud.   & se levanta      & he/she gets up \\
    nosotros/as       & nos levantamos  & we get up \\
    vosotros/as       & os levantáis    & you (pl.) get up \\
    ellos/as / Uds.   & se levantan     & they get up \\
    \bottomrule
\end{tabular}\end{center}

\begin{note}
    Recuerda a los estudiantes que el pronombre reflexivo va \emph{antes} del
    verbo conjugado (\textit{me levanto}), pero se une al infinitivo
    (\textit{antes de levantar\textbf{me}}). Es el error más común en esta unidad.
\end{note}

\section{Ejercicios}

\subsection{A. Pronombre reflexivo}
Completa con el pronombre reflexivo correcto.
\begin{enumerate}[leftmargin=*, label=\numbox{\arabic*}]
    \item Yo \fitb[1.4cm]{me} despierto a las siete.
    \item María \fitb[1.4cm]{se} ducha por la mañana.
    \item Nosotros \fitb[1.4cm]{nos} acostamos tarde.
    \item ¿Tú \fitb[1.4cm]{te} cepillas los dientes?
\end{enumerate}

\subsection{B. Conjugación}
Conjuga \textit{vestirse} para el sujeto indicado.
\begin{enumerate}[leftmargin=*, label=\numbox{\arabic*}, resume]
    \item
        \begin{exercise}[1.2cm]
            yo \rule{4cm}{0.4pt}
        \end{exercise}
        \begin{solution}
            yo me visto
        \end{solution}
    \item
        \begin{exercise}[1.2cm]
            ellos \rule{4cm}{0.4pt}
        \end{exercise}
        \begin{solution}
            ellos se visten
        \end{solution}
\end{enumerate}

\subsection{C. Traducción}
\begin{enumerate}[leftmargin=*, label=\numbox{\arabic*}, resume]
    \item
        \begin{exercise}[2.5cm]
            ``We wake up early and then we shower.''
        \end{exercise}
        \begin{solution}
            Nos despertamos temprano y luego nos duchamos.
        \end{solution}
\end{enumerate}

\subsection{D. Comprensión}
Responde según la lectura.
\begin{enumerate}[leftmargin=*, label=\numbox{\arabic*}, resume]
    \item
        \begin{exercise}[2cm]
            ¿A qué hora se levanta Lucía?
        \end{exercise}
        \begin{solution}
            Se levanta a las seis y media de la mañana.
        \end{solution}
    \item
        \begin{exercise}[2.5cm]
            ¿Por qué tiene sueño su hermano?
        \end{exercise}
        \begin{solution}
            Porque se acuesta muy tarde.
        \end{solution}
\end{enumerate}

\end{document}
